% ============================ CONCLUSIONES ===========================
\section{Conclusiones}

\paragraph{1. Cada base de datos resolvió la pregunta para la que fue elegida.} MongoDB conservó los 31.2 millones de eventos con su \texttt{payload} heterogéneo y permitió calcular el único indicador que depende de un campo anidado. Cassandra sirve los gráficos temporales desde tablas que suman menos de 1~MB, leyendo una partición por consulta. Neo4j responde preguntas de conectividad sobre 3.3 millones de nodos y 2.8 millones de relaciones. La justificación no es teórica: el KPI~4 no puede obtenerse de Cassandra, y el análisis del repositorio más conectado no puede obtenerse sin los tipos de relación del grafo.

\paragraph{2. La arquitectura por capas hizo que dos de las tres bases fueran reconstruibles.} Al alimentar Cassandra desde MongoDB y Neo4j desde Cassandra, y no las tres desde el archivo de origen, MongoDB queda como única fuente de la verdad. El grafo puede regenerarse completo sin volver a descargar ni a limpiar un solo archivo, y los indicadores se recalcularon para todo el periodo cuando cambió la definición del KPI~2. El costo es una dependencia secuencial: un fallo en una capa detiene las siguientes.

\paragraph{3. El \textit{pipeline} no pierde datos, y se comprobó en lugar de suponerlo.} Una auditoría independiente sobre los 985 archivos de origen contó 31,244,816 eventos crudos, de los cuales 190 se rechazan por validación y 402 por duplicados. Los 31,244,224 restantes coinciden exactamente con el número de documentos en MongoDB. Ninguna hora alcanzó el tope de eventos y ningún archivo contenía JSON inválido.

\paragraph{4. El principal hallazgo del proyecto es sobre la fuente, no sobre GitHub.} El volumen de los días completos varía en un factor de 25.6, pero esa variación no describe la actividad de la plataforma. Los eventos \texttt{PushEvent}, \texttt{CreateEvent} y \texttt{DeleteEvent}, que suman el 85.9\,\% del dataset, están presentes en 505 archivos horarios y ausentes en los otros 480, siempre los tres a la vez; el resto de los tipos se mantiene estable en todo el periodo. Los días con ese flujo completo promedian 1.74 millones de eventos y los días sin él, 108 mil. Un tablero que presentara el volumen diario como \guillemotleft{}tendencia de actividad\guillemotright{} estaría midiendo la cobertura del colector de GH Archive.

\paragraph{5. Un indicador solo es útil si se conoce qué mide.} La consecuencia práctica del punto anterior es que los KPI~1 y~2 no son comparables entre días con distinta cobertura, mientras que el KPI~4, apoyado en el flujo estable, se mantuvo entre 88.8\,\% y 94.7\,\% durante las seis semanas. De forma análoga, el promedio de 4.07 eventos por usuario convive con una mediana de 1, porque el 1\,\% de las cuentas genera el 43.4\,\% de la actividad y las cuentas automatizadas, el 21.7\,\%. Por eso el \textit{dashboard} acompaña cada gráfico con su lectura y declara sus limitaciones.

\paragraph{6. Un grado alto en el grafo no equivale a una comunidad.} El repositorio más conectado reúne 6,720 usuarios, cuatro veces más que el segundo, pero el 82\,\% de ellos solo le dio una estrella y apenas 7 cuentas hicieron \textit{push}. El modelo de grafo, al conservar el tipo de cada relación, permitió distinguir popularidad de colaboración; un simple conteo de usuarios por repositorio no lo habría hecho.

\paragraph{7. La idempotencia fue el requisito de ingeniería más exigente.} Cada base la resuelve de forma distinta: clave natural con \textit{upsert} en MongoDB, registro de horas procesadas en Cassandra, y la combinación de \texttt{MERGE} con ese mismo registro en Neo4j, porque \texttt{MERGE} evita duplicar aristas pero no evita sumar dos veces un contador. La limitación que subsiste está identificada: una escritura interrumpida a mitad de una hora puede acumular dos veces en las tablas diarias de Cassandra, y su remedio es recargar el día desde MongoDB.

\paragraph{8. Los fallos de operación no vinieron de los datos.} La única corrida fallida de Airflow se debió a un conflicto de permisos entre dos contenedores que comparten una carpeta, y se corrigió con una escritura atómica. Las dos horas que faltan en el registro corresponden a un intervalo en que el planificador no estuvo activo, que \texttt{catchup=False} no repone. En ambos casos el diagnóstico salió de los registros y de la tabla de trazabilidad \texttt{pipeline\_runs}, no del estado final de la corrida, la misma lección que dejó la primera entrega con los contadores de MapReduce.

\paragraph{9. Recomendaciones.} (a) Para análisis de volumen sobre este periodo conviene excluir el flujo de referencias de Git y trabajar con el flujo de colaboración, que es comparable entre días; el reporte público consultado menciona además un archivo alternativo compatible, mantenido por OpenDigger, que podría evaluarse como fuente. (b) El KPI~5 debería filtrarse por la propiedad \texttt{first\_seen} para depender de la fecha, como los otros cuatro. (c) Activar \texttt{catchup} acotado, o un sensor de horas faltantes sobre \texttt{pipeline\_runs}, haría que el sistema recupere por sí solo los intervalos en que el planificador se detiene. (d) Sustituir la acumulación por lectura y reescritura en Cassandra por columnas de tipo contador, o por el recálculo del día desde MongoDB, eliminaría la limitación de idempotencia descrita.
